\subsection{Métodos de descenso de primer orden}

Los métodos de descenso de primer orden construyen la dirección de búsqueda utilizando exclusivamente información de primer orden de la función objetivo, es decir, su gradiente.

La idea fundamental consiste en determinar una dirección que produzca una disminución local de la función al desplazarse desde el punto actual \(x_k\). Para analizar matemáticamente cómo varía la función ante pequeños desplazamientos, se introduce el concepto de derivada direccional.

Sea:
\[
p\in\mathbb{R}^n
\]

una dirección arbitraria. La derivada direccional de una función diferenciable \(f\) en el punto \(x\) y en la dirección \(p\) se define como:
\[
D_pf(x)=
\lim_{h\to0}
\frac{f(x+hp)-f(x)}{h}
\]

Esta expresión representa la tasa de variación instantánea de la función al desplazarse desde \(x\) en dirección \(p\).

Si la función es diferenciable, la derivada direccional puede escribirse como:
\[
D_pf(x)=\nabla f(x)^Tp
\]

La expresión:
\[
\nabla f(x)^Tp
\]

corresponde al producto interno entre el gradiente y el vector dirección \(p\). En consecuencia, la derivada direccional mide cuánto cambia localmente la función al proyectar el gradiente sobre la dirección considerada.

Para que un desplazamiento produzca una disminución local de la función debe cumplirse:
\[
D_pf(x)<0
\]

o equivalentemente:
\[
\nabla f(x)^Tp<0
\]

Toda dirección que satisfaga esta condición recibe el nombre de dirección de descenso.

Entre todas las posibles direcciones de descenso, resulta particularmente importante aquella que produce la máxima disminución local de la función. Para determinarla, se considera el siguiente problema:
\[
\min_{\|p\|=1}\nabla f(x)^Tp
\]

La restricción:
\[
\|p\|=1
\]

se introduce para evitar que el producto interno pueda hacerse arbitrariamente pequeño simplemente aumentando la magnitud del vector \(p\).

Aplicando la desigualdad de Cauchy-Schwarz:
\[
|\nabla f(x)^Tp|
\leq
\|\nabla f(x)\|\,\|p\|
\]

y utilizando la condición \(\|p\|=1\), se obtiene:
\[
|\nabla f(x)^Tp|
\leq
\|\nabla f(x)\|
\]

Por lo tanto:
\[
-\|\nabla f(x)\|
\leq
\nabla f(x)^Tp
\leq
\|\nabla f(x)\|
\]

El menor valor posible del producto interno corresponde entonces a:
\[
\nabla f(x)^Tp=
-\|\nabla f(x)\|
\]

La igualdad en la desigualdad de Cauchy-Schwarz se alcanza únicamente cuando ambos vectores poseen la misma dirección o direcciones opuestas. En este caso, el mínimo se obtiene cuando \(p\) posee dirección opuesta al gradiente. En consecuencia, la dirección de máximo descenso local viene dada por:
\[
p_k=-\nabla f(x_k)
\]

Esta elección conduce al método de descenso por gradiente o \textit{steepest descent}, cuya actualización iterativa toma la forma:
\[
x_{k+1}=x_k-\alpha_k\nabla f(x_k)
\]

Una vez determinada la dirección de descenso, resta establecer cuánto avanzar sobre dicha dirección. El parámetro \(\alpha_k\) controla precisamente la magnitud del desplazamiento realizado durante cada iteración.

Para determinar el valor adecuado de \(\alpha_k\), se restringe la función objetivo a la recta generada por la dirección de búsqueda. Esto se realiza definiendo la función auxiliar:
\[
\phi(\alpha)=f(x_k+\alpha p_k)
\]

La variable \(\alpha\) parametriza todos los puntos de la recta que pasa por \(x_k\) en dirección \(p_k\). De esta manera, el problema multivariable original se reduce localmente a un problema de optimización unidimensional.

El objetivo de la búsqueda lineal (\textit{line search}) consiste en encontrar un valor de \(\alpha_k\) tal que:
\[
\phi(\alpha_k)<\phi(0)
\]

garantizando así una disminución efectiva de la función objetivo en cada iteración.

A pesar de su simplicidad conceptual y bajo costo computacional, los métodos de descenso de primer orden presentan limitaciones importantes. Cuando la función posee curvaturas muy diferentes según la dirección considerada, el algoritmo puede exhibir trayectorias oscilatorias y convergencia lenta. Este comportamiento surge debido a que el método utiliza únicamente información de pendiente local, sin incorporar información explícita sobre la curvatura de la superficie de optimización.

Estas limitaciones motivan el desarrollo de métodos de segundo orden, capaces de incorporar información de curvatura local mediante derivadas segundas de la función objetivo.